\section{Related Work}
\label{sec:related}

Hybrid architectures combining convolutional and recurrent layers,
often with a further decomposition or gradient-boosted stage, dominate
recent work on short-term PV power forecasting, and reported accuracy
gains are large: a meta-analysis of 180 studies published since 2007
concludes that hybrid models consistently outperform alternatives
\cite{nguyenmusgens2021drives}. That conclusion aggregates reported
errors; this paper asks what those errors are measured against.

To characterise current practice, we coded 27 papers on hybrid PV
power forecasting published between 2022 and 2026, on eight dimensions
of evaluation protocol, each judgement supported by a verbatim
quotation from the source (Table~\ref{tab:T1}). Beyond the two counts
given in Section~\ref{sec:intro}, seventeen of 27 papers compare only
against their own architecture's components, and seventeen do not
state whether night hours are excluded from evaluation - the two
protocol axes this paper measures.

\input{tables/T1_survey.tex}

Rigorous verification practice exists in solar forecasting. Yang et
al.~\cite{yang2020verification}, writing with more than thirty
co-authors, recommend universally reporting the skill score against
the optimal convex combination of climatology and persistence. One
paper in our sample of 27 meets that standard: Mayer
\cite{mayer2022physmlhybrid} evaluates day-ahead forecasts for
fourteen plants against that convex reference, states its daylight
filter as a protocol, and holds out a full year used in no part of
model selection; its headline hybrid improvement is 5.2 percent in
mean absolute error and 1.0 percent in root-mean-square error over an
optimised physical model chain, an order of magnitude smaller than
improvements typically reported elsewhere in the sample, and measured
against a considerably stronger baseline. The standard is achievable
and is published in the same journals the surveyed papers appear in;
we adopt its reference forecast here and hold the model fixed while
varying it.

The difficulty of comparing results across studies with differing
evaluation protocols is not specific to solar forecasting; it has been
documented across at least 294 papers in 17 scientific fields,
together with a taxonomy of the mechanisms involved
\cite{kapoor2023leakage}.
